% A LaTeX template for MSc Thesis submissions to 
% Politecnico di Milano (PoliMi) - School of Industrial and Information Engineering
%
% S. Bonetti, A. Gruttadauria, G. Mescolini, A. Zingaro
% e-mail: template-tesi-ingind@polimi.it
%
% Last Revision: October 2021
%
% Copyright 2021 Politecnico di Milano, Italy. NC-BY

\documentclass{Configuration_Files/PoliMi3i_thesis}

%------------------------------------------------------------------------------
%	REQUIRED PACKAGES AND  CONFIGURATIONS
%------------------------------------------------------------------------------

% CONFIGURATIONS
\usepackage{parskip} % For paragraph layout
\usepackage{setspace} % For using single or double spacing
\usepackage{emptypage} % To insert empty pages
\usepackage{multicol} % To write in multiple columns (executive summary)
\setlength\columnsep{15pt} % Column separation in executive summary
\setlength\parindent{0pt} % Indentation
\raggedbottom  

% PACKAGES FOR TITLES
\usepackage{titlesec}
% \titlespacing{\section}{left spacing}{before spacing}{after spacing}
\titlespacing{\section}{0pt}{3.3ex}{2ex}
\titlespacing{\subsection}{0pt}{3.3ex}{1.65ex}
\titlespacing{\subsubsection}{0pt}{3.3ex}{1ex}
\usepackage{color}

% PACKAGES FOR LANGUAGE AND FONT
\usepackage[english]{babel} % The document is in English  
\usepackage[utf8]{inputenc} % UTF8 encoding
\usepackage[T1]{fontenc} % Font encoding
\usepackage[11pt]{moresize} % Big fonts

% PACKAGES FOR IMAGES
\usepackage{graphicx}
\usepackage{transparent} % Enables transparent images
\usepackage{eso-pic} % For the background picture on the title page
\usepackage{subfig} % Numbered and caption subfigures using \subfloat.
\usepackage{tikz} % A package for high-quality hand-made figures.
\usetikzlibrary{}
\graphicspath{{./Images/}} % Directory of the images
\usepackage{caption} % Coloured captions
\usepackage{xcolor} % Coloured captions
\usepackage{amsthm,thmtools,xcolor} % Coloured "Theorem"
\usepackage{float}

% STANDARD MATH PACKAGES
\usepackage{amsmath}
\usepackage{amsthm}
\usepackage{amssymb}
\usepackage{amsfonts}
\usepackage{bm}
\usepackage[overload]{empheq} % For braced-style systems of equations.
\usepackage{fix-cm} % To override original LaTeX restrictions on sizes

% PACKAGES FOR TABLES
\usepackage{tabularx}
\usepackage{longtable} % Tables that can span several pages
\usepackage{colortbl}

% PACKAGES FOR ALGORITHMS (PSEUDO-CODE)
\usepackage{algorithm}
\usepackage{algorithmic}

% PACKAGES FOR REFERENCES & BIBLIOGRAPHY
\usepackage[colorlinks=true,linkcolor=black,anchorcolor=black,citecolor=black,filecolor=black,menucolor=black,runcolor=black,urlcolor=black]{hyperref} % Adds clickable links at references
\usepackage{cleveref}
\usepackage[square, numbers, sort&compress]{natbib} % Square brackets, citing references with numbers, citations sorted by appearance in the text and compressed
\bibliographystyle{abbrvnat} % You may use a different style adapted to your field

% OTHER PACKAGES
\usepackage{pdfpages} % To include a pdf file
\usepackage{afterpage}
\usepackage{lipsum} % DUMMY PACKAGE
\usepackage{fancyhdr} % For the headers
\fancyhf{}

% Input of configuration file. Do not change config.tex file unless you really know what you are doing. 
% Define blue color typical of polimi
\definecolor{bluepoli}{cmyk}{0.4,0.1,0,0.4}

% Custom theorem environments
\declaretheoremstyle[
  headfont=\color{bluepoli}\normalfont\bfseries,
  bodyfont=\color{black}\normalfont\itshape,
]{colored}

% Set-up caption colors
\captionsetup[figure]{labelfont={color=bluepoli}} % Set colour of the captions
\captionsetup[table]{labelfont={color=bluepoli}} % Set colour of the captions
\captionsetup[algorithm]{labelfont={color=bluepoli}} % Set colour of the captions

\theoremstyle{colored}
\newtheorem{theorem}{Theorem}[chapter]
\newtheorem{proposition}{Proposition}[chapter]

% Enhances the features of the standard "table" and "tabular" environments.
\newcommand\T{\rule{0pt}{2.6ex}}
\newcommand\B{\rule[-1.2ex]{0pt}{0pt}}

% Pseudo-code algorithm descriptions.
\newcounter{algsubstate}
\renewcommand{\thealgsubstate}{\alph{algsubstate}}
\newenvironment{algsubstates}
  {\setcounter{algsubstate}{0}%
   \renewcommand{\STATE}{%
     \stepcounter{algsubstate}%
     \Statex {\small\thealgsubstate:}\space}}
  {}

% New font size
\newcommand\numfontsize{\@setfontsize\Huge{200}{60}}

% Title format: chapter
\titleformat{\chapter}[hang]{
\fontsize{50}{20}\selectfont\bfseries\filright}{\textcolor{bluepoli} \thechapter\hsp\hspace{2mm}\textcolor{bluepoli}{|   }\hsp}{0pt}{\huge\bfseries \textcolor{bluepoli}
}

% Title format: section
\titleformat{\section}
{\color{bluepoli}\normalfont\Large\bfseries}
{\color{bluepoli}\thesection.}{1em}{}

% Title format: subsection
\titleformat{\subsection}
{\color{bluepoli}\normalfont\large\bfseries}
{\color{bluepoli}\thesubsection.}{1em}{}

% Title format: subsubsection
\titleformat{\subsubsection}
{\color{bluepoli}\normalfont\large\bfseries}
{\color{bluepoli}\thesubsubsection.}{1em}{}

% Shortening for setting no horizontal-spacing
\newcommand{\hsp}{\hspace{0pt}}

\makeatletter
% Renewcommand: cleardoublepage including the background pic
\renewcommand*\cleardoublepage{%
  \clearpage\if@twoside\ifodd\c@page\else
  \null
  \AddToShipoutPicture*{\BackgroundPic}
  \thispagestyle{empty}%
  \newpage
  \if@twocolumn\hbox{}\newpage\fi\fi\fi}
\makeatother

%For correctly numbering algorithms
\numberwithin{algorithm}{chapter}

%----------------------------------------------------------------------------
%	NEW COMMANDS DEFINED
%----------------------------------------------------------------------------

% EXAMPLES OF NEW COMMANDS
\newcommand{\bea}{\begin{eqnarray}} % Shortcut for equation arrays
\newcommand{\eea}{\end{eqnarray}}
\newcommand{\e}[1]{\times 10^{#1}}  % Powers of 10 notation

%----------------------------------------------------------------------------
%	ADD YOUR PACKAGES (be careful of package interaction)
%----------------------------------------------------------------------------

%----------------------------------------------------------------------------
%	ADD YOUR DEFINITIONS AND COMMANDS (be careful of existing commands)
%----------------------------------------------------------------------------
\newtheorem{definition}{Definition}
\newcommand{\mc}[1]{\mathcal{#1}}
\DeclareMathOperator*{\argmax}{arg\,max}
\DeclareMathOperator*{\argmin}{arg\,min}
%----------------------------------------------------------------------------
%	BEGIN OF YOUR DOCUMENT
%----------------------------------------------------------------------------
\includeonly{chapters/Preliminaries}
\begin{document}

\fancypagestyle{plain}{%
\fancyhf{} % Clear all header and footer fields
\fancyhead[RO,RE]{\thepage} %RO=right odd, RE=right even
\renewcommand{\headrulewidth}{0pt}
\renewcommand{\footrulewidth}{0pt}}

%----------------------------------------------------------------------------
%	TITLE PAGE
%----------------------------------------------------------------------------

\pagestyle{empty} % No page numbers
\frontmatter % Use roman page numbering style (i, ii, iii, iv...) for the preamble pages

\puttitle{
	title=Title, % Title of the thesis
	name=Thomas Bonenfant, % Author Name and Surname
	course=Computer Science and Engineering - Ingegneria Informatica, % Study Programme (in Italian)
	ID  = 000000,  % Student ID number (numero di matricola)
	advisor= Prof. Marcello Restelli, % Supervisor name
	coadvisor={Amarildo Likmeta, Davide Salaorni}, % Co-Supervisor name, remove this line if there is none
	academicyear={2023-24},  % Academic Year
} % These info will be put into your Title page
%----------------------------------------------------------------------------
%	PREAMBLE PAGES: ABSTRACT (inglese e italiano), EXECUTIVE SUMMARY
%----------------------------------------------------------------------------
\startpreamble
\setcounter{page}{1} % Set page counter to 1

% ABSTRACT IN ENGLISH
\chapter*{Abstract} 
Air Hockey is a competitive, dynamic, and complex game. 
Utilizing robotic systems to play Air Hockey presents various challenges, including physical limitations and constraints that must 
be respected to ensure safety and prevent damage to the components. 
This thesis explores the application of Reinforcement Learning techniques in this context. 
Specifically, since Air Hockey comprises several sub-problems, such as hitting the puck to score a point or defending one's goal, 
a hierarchical optimization problem is formulated to develop an agent capable of playing entire matches while adhering to physical constraints. 
This agent is organized into two levels. The lower level comprises several specialized policies, trained using Deep RL and Rule-based RL algorithms, 
while the higher level selects which specialized policy to use through a parametric state machine, optimized using Rule-Based RL. 
Additionally, it describes our participation in the Robot Air Hockey Challenge, a competition designed to address the \textit{sim-to-real gap}, 
the performance and safety drop observed when a real robotic system is controlled by a policy trained in a simulated environment.
\\
\\
\textbf{Keywords:} Robotics, Air Hockey, Rule-based policies, Reinforcement Learning, Deep Reinforcement Learning, Hierarchical Reinforcement Learning,
Explainable Artificial Intelligence % Keywords

% ABSTRACT IN ITALIAN
\chapter*{Abstract in lingua italiana}
% Air Hockey è un gioco competitivo, dinamico e complesso. L'impiego di sistemi robotici per giocare ad Air Hockey presenta varie sfide, 
% tra cui limiti fisici e vincoli che devono essere rispettati per garantire la sicurezza ed evitare danni ai componenti. 
% Questa tesi esplora l'applicazione di tecniche di Reinforcement Learning in questo contesto. 
% In particolare, dato che Air Hockey è composto da diversi sottoproblemi, come colpire il disco per segnare un punto o difendere la propria porta, 
% viene formulato un problema di ottimizzazione gerarchica per sviluppare un agente capace di giocare intere partite rispettando i vincoli fisici. 
% Inoltre, viene descritta la nostra partecipazione alla Robot Air Hockey Challenge, una competizione ideata per affrontare il \textit{sim-to-real gap}, 
% ovvero il calo di performance e sicurezza che si verifica quando un sistema robotico reale è controllato da una politica allenata in un ambiente simulato.\\

Air Hockey è un gioco competitivo, dinamico e complesso. L'impiego di sistemi robotici per giocare ad Air Hockey presenta varie sfide, 
tra cui limiti fisici e vincoli che devono essere rispettati per garantire la sicurezza ed evitare danni ai componenti. 
Questa tesi esplora l'applicazione di tecniche di Reinforcement Learning in questo contesto. 
In particolare, dato che Air Hockey è composto da diversi sottoproblemi, come colpire il disco per segnare un punto o difendere la propria porta, 
viene formulato un problema di ottimizzazione gerarchica per sviluppare un agente capace di giocare intere partite rispettando i vincoli fisici.
Questo agente è organizzato in due livelli. Il livello più basso comprende diverse politiche specializzate, che vengono allenate con algoritmi di Deep RL
e Rule-based RL, mentre il livello superiore seleziona quale politica specializzata utilizzare tramite una macchina a stati parametrica, ottimizzata tramite Rule-Based RL.
Inoltre, viene descritta la nostra partecipazione alla Robot Air Hockey Challenge, una competizione ideata per affrontare il \textit{sim-to-real gap}, 
ovvero il calo di performance e sicurezza che si verifica quando un sistema robotico reale è controllato da una politica allenata in un ambiente simulato.\\

\textbf{Parole chiave:} Robotica, Air Hockey, Politiche a regole, Reinforcement Learning, Deep Reinforcement Learning, Hierarchical Reinforcement Learning, Explainable Artificial Intelligence % Keywords (italian)
%----------------------------------------------------------------------------
%	LIST OF CONTENTS/FIGURES/TABLES/SYMBOLS
%----------------------------------------------------------------------------

% TABLE OF CONTENTS
\thispagestyle{empty}
\tableofcontents % Table of contents 
\thispagestyle{empty}
\cleardoublepage

\addtocontents{toc}{\vspace{2em}} % Add a gap in the Contents, for aesthetics


\mainmatter % Begin numeric (1,2,3...) page numbering

% --------------------------------------------------------------------------
% NUMBERED CHAPTERS % Regular chapters following
% --------------------------------------------------------------------------
\chapter*{Introduction}

This document is intended to be both an example of the Polimi \LaTeX{} template for Master Theses,
as well as a short introduction to its use. It is not intended to be a general introduction to \LaTeX{} itself,
and the reader is assumed to be familiar with the basics of creating and compiling \LaTeX{} documents (see \cite{oetiker1995not, kottwitz2015latex}). 
\\
The cover page of the thesis must contain all the relevant information:
title of the thesis, name of the Study Programme and School, name of the author,
student ID number, name of the supervisor, name(s) of the co-supervisor(s) (if any), academic year.
The above information are provided by filling all the entries in the command \verb|\puttitle{}|
in the title page section of this template.
\\
Be sure to select a title that is meaningful.
It should contain important keywords to be identified by indexer.
Keep the title as concise as possible and comprehensible even to people who are not experts in your field.
The title has to be chosen at the end of your work so that it accurately captures the main subject of the manuscript. 
\\
Since a thesis might be a substantial document, it is convenient to break it into chapters.
You can create a new chapter as done in this template by simply using the following command
\begin{verbatim}
\chapter{Title of the chapter}
\end{verbatim}
followed by the body text.
\\
Especially for long manuscripts, it is recommended to give each chapter its own file.
In this case, you write your chapter in a separated \verb|chapter_n.tex| file
and then include it in the main file with the following command
\begin{verbatim}
\input{chapter_n.tex}
\end{verbatim}
It is recommended to give a label to each chapter by using the command
\begin{verbatim}
\label{ch:chapter_name}%
\end{verbatim}
where the argument is just a text string that you'll use to reference that part
as follows: \textit{Chapter~\ref{ch:chapter_one} contains \sc{an introduction to}  \dots}.\\
If necessary, an unnumbered chapter can be created by
\begin{verbatim}
\chapter*{Title of the unnumbered chapter}
\end{verbatim}

\chapter{Preliminaries}
\label{ch:Preliminaries}

This chapter starts introducing what a Markov Decision Process is and how it is used to formally define a sequential decision problem.
Next it describes what a policy is and its associated utility functions used to evaluate its performance.

\section{Markov Decision Process}
In a Sequential Decision problem an agent interacts with an environment through actions and receives a feedback signal from it.
The environment is usually modeled as a Markov Decision Process.
\\
\begin{definition}
A Finite Markov Decision Process (MDP) is a tuple $\mc{M} = (\mc{S}, \mc{A}, p, r, \gamma, \mu)$ specified by:
\begin{itemize}
    \item a finite set $\mc{S}$ of possible states
    \item a finite set $\mc{A}$ of actions
    \item a transition probability function $p: \mc{S} \times \mc{A} \times \mc{S} \to [0, 1]$
    \item a reward function $r: \mc{S} \times \mc{A} \to \mathbb{R} $
    \item a discount factor $\gamma \in (0,1)$
    \item a probability distribution over the initial states $\mu: \mc{S} \to [0, 1]$

\end{itemize}
\end{definition}

The agent interacts with it at discrete time steps $t = 0,1,2,3,\dots$.
At each time step the agent receives some representation of the environment's \textit{state}, $S_t \in \mc{S}$, and on that basis
selects an \textit{action}, $A_t \in \mc{A}(s)$. One time step later, the agent receives a numerical \textit{reward}, $R_{t+1} \in \mc{R} \subset \mathbb{R}$,
and finds itself in a new state, $S_{t+1}$. The MDP and agent together thereby give rise to a sequence or \textit{trajectory} that begins like this:
\[ S_0, A_0, R_1, S_1, A_1, R_2, S_2, A_2, R_3, \dots\]

The transition probability depends only on the previous state and action thanks to the Markov assumption:
\[p(s'|s,a) \doteq \mathbb{P}(S_t = s' | S_{t-1}=s, A_{t-1}=a)\]
The reward function is the expected reward the state action pair:
\[r(s,a) \doteq \mathbb{E}[R_t | S_{t-1} = s, A_{t-1} = a]\]

The agent choses actions according to a policy. A policy $\pi$ is a distribution over actions given the state:
\[\pi(a|s)=\mathbb{P}(a|s)\]
Therefore the policy can be stochastic or deterministic and because it considers only the previous state it is also Markovian.
The goal of the agent is to learn a policy $\pi$ that maximizes the expected return, where the return is defined as the cumulative discounted reward:
\[G_t \doteq R_{t+1} + \gamma R_{t+2} + \gamma^2 R_{t+3} + \dots = \sum_{k=0}^{\infty}\gamma^k R_{t+k+1}\]

Given a policy $\pi$, it is possible to define the utility of each state through a value function.
\begin{definition}
    The state value function $v^\pi: \mc{S} \to \mathbb{R}$ of an MDP is the expected return starting from state s and then following policy $\pi$
    \[v^\pi(s) \doteq \mathbb{E}_\pi[G_t | S_t=s]\]
\end{definition}

$\mathbb{E}_\pi[.]$ denotes the expected value of a random variable given that the agent follows policy $\pi$. 
For control purposes, rather than the value of each state we consider the value of each action in each state.
\begin{definition}
    The action-value function $q^\pi: \mc{S} \times \mc{A} \to \mathbb{R}$ of an MDP is the expected return starting from state s, taking action a and then following policy $\pi$
    \[q^\pi(s,a) \doteq \mathbb{E}_\pi[G_t | S_t=s, A_t=a]\]
\end{definition}

A fundamental property of the value functions is that they satisfy
the following recursive relationships, called \textit{Bellman expectation equations}:

\begin{equation}
    \label{eq:bellman_v}
    v^\pi(s) = \sum_{a \in \mc{A}(s)}\pi(a | s)\left( r(s, a) + \gamma \sum_{s' \in \mc{S}} p(s'|s,a) v^\pi(s') \right)
\end{equation}

\begin{equation}
\label{eq:bellman_q}
q^\pi(s,a) = r(s,a) + \gamma \sum_{s' \in \mc{S}} p(s'|s,a)\sum_{a' \in \mc{A}(s')}\pi(a'|s')q^\pi(s',a')
\end{equation}

% DIRECT SOLUTION FOR V

\begin{definition}
    The Markov Reward Process induced by a policy $\pi$ on a MDP is the tuple $(\mc{S}, p^\pi, r^\pi)$ specified by:
    \begin{itemize}
    \item the same set $\mc{S}$ from the MDP
    \item the transition probability function $p^\pi: \mc{S} \times \mc{S} \to [0,1]$ obtained from following policy $\pi$
    \item the reward function $r: \mc{S} \to \mathbb{R}$ obtained fom following policy $\pi$
    \end{itemize}
\end{definition}
\begin{align*}
    p^\pi(s'|s) &= \mathbb{P}_\pi(S_t=s'|S_{t-1}=s) = \sum_{a \in \mc{A}(s)}\pi(a|s) p(s'|s,a) \\
    r^\pi(s) &= \sum_{a \in \mc{A}(s)}\pi(a|s) r(s, a)
\end{align*}

The bellman expectation equation for $v^\pi$ can be written in matrix form using the induced MRP and a closed form solution can be computed:

\begin{align*}
    V^\pi &= R^\pi + \gamma P^\pi V^\pi \\
    V^\pi &= (I - P^\pi)^{-1} R^\pi
\end{align*}

where $V^\pi \in \mathbb{R}^{|\mc{S}|}$ is the matrix equivalent of $v^\pi$.

% BELLMAN OPERATOR

\begin{definition}
    The Bellman operator for $V^\pi$ is defined as $T^\pi : \mathbb{R}^{|\mc{S}|} \to \mathbb{R}^{|\mc{S}|}$
    \[(T^\pi v^\pi)(s) = \sum_{a \in \mc{A}(s)}\pi(a|s)\left(r(s,a) + \gamma \sum_{s' \in \mc{S}}p(s'|s,a)v^\pi(s')\right)\]
\end{definition}

Using the Bellman operator for $v\pi$, the Bellman expectation equation can be compactly written as:
\begin{equation*}
    T^\pi V^\pi = V^\pi
\end{equation*}

% The Bellman operator for $v^\pi$ exhibits the following properties:
% \begin{enumerate}
%     \item $v^\pi$ is a fixed point of the Bellman operator $T^\pi$
%     \item The Bellman expectation equation is linear in $v^\pi$ and $T^\pi$
%     \item If $0 < \gamma < 1$ then $T^\pi$ is a contraction with respect to the maximum norm
% \end{enumerate}

\begin{definition}
    The Bellman operator for $q^\pi$ is defined as $T^\pi : \mathbb{R}^{|\mc{S} \times \mc{A}|} \to \mathbb{R}^{|\mc{S} \times \mc{A}|}$
    \[(T^\pi q^\pi)(s, a) = r(s,a) + \gamma \sum_{s' \in \mc{S}}p(s'|s,a)\sum_{a' \in \mc{A}(s')}\pi(a'|s')q^\pi(s,a)\]
\end{definition}

Using the Bellman operator for $q^\pi$, the Bellman expectation equations can be compactly written as:
\begin{equation*}
    T^\pi Q^\pi = Q^\pi
\end{equation*}

% The Bellman operator for $q^\pi$ exhibits the following properties:
% \begin{enumerate}
%     \item $q^\pi$ is a fixed point of the Bellman operator $T^\pi$
%     \item The Bellman expectation equation is linear in $q^\pi$ and $T^\pi$
%     \item If $0 < \gamma < 1$ then $T^\pi$ is a contraction with respect to the maximum norm
% \end{enumerate}

% OPTIMAL VALUE FUNCTIONS

\begin{definition}
    The optimal state-value function $v^*(s)$ is the maximum state-value function over all policies
    \[v^*(s) = \max_\pi v^\pi(s)\]
\end{definition}
\begin{definition}
    The optimal action-value function $q^*(s)$ is the maximum action-value function over all policies
    \[q^*(s,a) = \max_\pi q^\pi(s,a)\]
\end{definition}

Also the optimal value functions satisfy a recursive relationship, called \textit{Bellman Optimality Equations}

\begin{equation}
    \label{eq:bellman_optimality_v}
    \begin{aligned}
        v^*(s) &= \max_{a \in \mc{A}(s)} q^*(s,a) \\
        &= \max_{a \in \mc{A}(s)} \left(r(s,a) + \gamma \sum_{s' \in \mc{S}}p(s'|s,a)v^*(s)\right)
    \end{aligned}
\end{equation}

\begin{equation}
    \label{eq:bellman_optimality_q}
    \begin{aligned}
        q^*(s, a) &= r(s,a) + \gamma \sum_{s' \in \mc{S}}p(s'|s,a)v^*(s')  \\
        &= r(s,a) + \gamma \sum_{s' \in \mc{S}}p(s'|s,a)\max_{a' \in \mc{A}(s)} q^*(s',a')
    \end{aligned}
\end{equation}

Value functions define a partial ordering over policies
\[\pi > \pi' \text{ if } v^\pi(s) > v^{\pi'}(s),\quad \forall s \in \mc{S}\]

\begin{theorem}
    For any Markov Decision Process
    \begin{enumerate}
        \item There exists an optimal policy $\pi*$ that is better than or equal to all other policies \\
        $\pi^* \ge \pi, \quad \forall \pi$
        \item All optimal policies achieve the optimal value function $v^{\pi^*}(s) = v*(s)$
        \item All optimal policies achieve the optimal state-value function $q^{\pi^*}(s, a) = q^*(s,a)$
        \item There is always a deterministic optimal policy for an MDP
    \end{enumerate}
\end{theorem}

A determinisitc optimal policy can be found by maximizing over $q^*(s,a)$ (greedy policy)
\begin{equation*}
    \pi^*(s,a) = 
    \begin{cases}
        1 \quad & \text{if } a = \argmax_{a \in \mc{A}}q^*(s,a)\\
        0 \quad & \text{otherwise}
    \end{cases}
\end{equation*}

% OPTIMAL BELLMAN OPERATOR

\begin{definition}
    The Bellman optimality operator for $v^*$ is defined as $T^* : \mathbb{R}^{|\mc{S}|} \to \mathbb{R}^{|\mc{S}|}$
    \[(T^* v^*)(s) = \max_{a \in \mc{A}(s)}\left(r(s,a) + \gamma \sum_{s' \in \mc{S}}p(s'|s,a)v^*(s')\right)\]
\end{definition}

\begin{definition}
    The Bellman optimality operator for $q^*$ is defined as $T^* : \mathbb{R}^{|\mc{S} \times \mc{A}|} \to \mathbb{R}^{|\mc{S} \times \mc{A}|}$
    \[(T^* q^*)(s, a) = r(s,a) + \gamma \sum_{s' \in \mc{S}}p(s'|s,a)\max_{a' \in \mc{A}(s')}q^*(s',a')\]
\end{definition}

The Bellman operators exhibit the following properties:
\begin{itemize}
    \item \textbf{Monotonicity}: if $f_1 \le f_2$ component-wise \\
    \[T^\pi f_1 \le T^\pi f_2 \quad,\quad T^*f_1 \le T^*f_2\]
    \item \textbf{Max-Norm Contraction}: for two vectors $f_1$ and $f_2$ \\
    \begin{align*}
    \lVert T^\pi f_1 - T^\pi f_2 \rVert _\infty &\le \gamma \lVert f_1 - f_2 \rVert _\infty \\
    \lVert T^* f_1 - T^* f_2 \rVert _\infty &\le \gamma \lVert f_1 - f_2 \rVert _\infty
    \end{align*}
    \item $v^\pi$ is the \textbf{unique fixed point} of $T^\pi$
    \item $v^*$ is the \textbf{unique fixed point} of $T^*$
    \item For any vector $f \in \mathbb{R}^{|\mc{S}|}$ and any policy $\pi$, we have \\
    \[\lim_{k \to \infty}(T^\pi)^kf = v^\pi \quad,\quad \lim_{k \to \infty}(T^*)^kf = v^*\]
    
\end{itemize}

Differently from \eqref{eq:bellman_v}, the Bellman optimality equation is non linear due to the presence of a max operation
and no closed form solution exists for the general case.
\\
The bellman optimality equation of an MDP is solved by different iterative methods like Dynamic programming, Linear programming and Reinforcement Learning.


\section{Dynamic Programming \label{sec:dynamic_programming}}
The term dynamic programming (DP) refers to a collection of algorithms that can be used to compute optimal policies given a perfect model
of the environment as a Markov Decision Process. Classical DP algorithms are of limited utility in reinforcement learning both because of their assumption
of a perfect model and because of their computational expense

\subsection{Policy Evaluation (Prediction)}
In Dynamic Programming, policy evaluation is done iteratively using the bellman equation for $v^\pi(s)$ \eqref{eq:bellman_v} as an update rule.
This update rule is called Bellman Expectation Backup and applying it to each state is called a \textit{sweep}.
\begin{equation*}
    \label{eq:policy_evaluation_update}
    \begin{aligned}
    v_{k+1}(s) &\doteq \mathbb{E}_\pi[R_{t+1} + \gamma v_k^{S_{t+1}} | S_t = s] \\
    &= \sum_{a \in \mc{A}} \pi(a|s)[r(s,a) +\gamma \sum_{s' \in \mc{S}} p(s'|s,a) v_k(s')]
    \end{aligned}
\end{equation*}

\begin{algorithm}[H]
    \label{alg:iterative_policy_evaluation}
    \caption{Iterative Policy Evaluation}
    \begin{algorithmic}[1]
        \STATE \textbf{Input:} A Markov Decision Process (MDP) defined by $(S, A, P, R, \gamma)$ and a policy $\pi$
        \STATE \textbf{Output:} State-value function $V$ for policy $\pi$
        \STATE \textbf{Initialize} $V(s) \gets 0$ for all $s \in S$
        \STATE \textbf{Initialize} $\theta \gets$ a small positive number (convergence threshold)
        \REPEAT
            \STATE $\Delta \gets 0$
            \FOR{each $s \in S$}
                \STATE $v \gets V(s)$
                \STATE $V(s) \gets \sum_{a \in \mc{A}} \pi(a|s) [r(s,a) + \gamma \sum_{s' \in \mc{S}} P(s'| s, a) V(s')]$
                \STATE $\Delta \gets \max(\Delta, |v - V(s)|)$
            \ENDFOR
        \UNTIL{$\Delta < \theta$}
        \STATE \textbf{return} $V$
    \end{algorithmic}
\end{algorithm} 

\subsection{Policy Improvement}
Considering a deterministic policy $\pi$ we can improve it by acting greedily:
\[\pi'(s) = \argmax_{a \in \mc{A}}q^\pi(s,a)\]

\begin{theorem}
    \label{th:policy_improvement}
    Let $\pi$ and $\pi'$ be any pair of deterministic policies such that
    \begin{equation*}
    q^\pi(s,\pi'(s)) \ge v^\pi(s) \quad, \quad \forall s \in \mc{S}
    \end{equation*}
    Then the policy $\pi'$ must be as good as, or better than $\pi$
    \begin{equation}
    v^{\pi'}(s) \ge v^\pi(s) \quad, \quad \forall s \in \mc{S}
    \end{equation}
\end{theorem}

If the improvement stops ($v^{\pi'}(s) = v^\pi(s) \, \forall s \in \mc{S}$) then the Bellman Optimality Equation \eqref{eq:bellman_optimality_v}
holds and $\pi' = \pi$ is an optimal policy $\pi^*$.

The policy iteration algorithm \ref{alg:policy_iteration} alternates between policy evaluation and policy improvements until the policy is stable and
thus optimal.

\begin{algorithm}[H]
    \caption{Policy Iteration}
    \label{alg:policy_iteration}
    \begin{algorithmic}[1]
        \STATE \textbf{Input:} A Markov Decision Process (MDP) defined by $(\mc{S}, \mc{A}, p, r, \gamma)$ and a policy $\pi$
        \STATE \textbf{Output:} Optimal policy $\pi^*$ and Optimal state-value function $V^*(s)$
        \STATE \textbf{Initialize} $V(s) \in \mathbb{R}$ and $\pi(s) \in \mc{A}(s)$ arbitrarily for all $s \in \mc{S}$
        \STATE \textbf{Initialize} $\theta \gets$ a small positive number (convergence threshold)
        \REPEAT
            \STATE $\Delta \gets 0$
            \FOR{each $s \in \mc{S}$}
                \STATE $v \gets V(s)$
                \STATE $V(s) \gets \sum_{a \in \mc{A}} \pi(a|s) [r(s,a) + \gamma \sum_{s' \in \mc{S}} P(s' | s, a) V(s')]$
                \STATE $\Delta \gets \max(\Delta, |v - V(s)|)$
            \ENDFOR
        \UNTIL{$\Delta < \theta$}
        
        \STATE \textit{policy-stable} $\gets$ \textit{true}
        \FOR{each $s \in \mc{S}$}
            \STATE \textit{old-action} $\gets \pi(s)$
            \STATE $\pi(s) \gets \argmax_{a \in \mc{A}(s)}[r(s,a) + \gamma \sum_{s' \in \mc{S}}p(s'|s,a) * V(s)]$
            \IF{\textit{old-action} $\neq \pi(s)$}
                \STATE \textit{policy-stable} $\gets$ \textit{false}
            \ENDIF
        \ENDFOR
        \IF{\textit{policy-stable}}
            \STATE \textbf{return} $V \approx v^*$ and $\pi \approx \pi^*$
        \ELSE
            \STATE go to 5
        \ENDIF
    \end{algorithmic}
\end{algorithm}

A small modification to the policy iteration algorithms lead to the Value iteration algorithm \ref{alg:value_iteration} in which the evaluation is stopped after only one sweep
and followed by the policy improvement. It basically implements the bellman optimality equation \eqref{eq:bellman_optimality_v} as an update rule:
\begin{equation*}
    v_{k+1}(s) = \max_{a \in \mc{A}}\left(r(s,a) + \gamma\sum_{s' \in \mc{S}}p(s'|s,a)v_k(s)\right)
\end{equation*}

Once $v_k \approx v^*$ we can build a greedy policy $\pi^*$

\begin{algorithm}[H]
    \caption{Value Iteration}\label{alg:value_iteration}
    \begin{algorithmic}[1]
    \STATE Initialize $V(s) = 0$ for all states $s \in S$
    \REPEAT
        \STATE $\Delta \gets 0$
        \FOR{$s \in S$}
            \STATE $v \gets V(s)$
            \STATE $V(s) \gets \max_{a \in A(s)} \sum_{s'} P(s'|s,a) \left[ R(s,a) + \gamma V(s') \right]$
            \STATE $\Delta \gets \max(\Delta, |v - V(s)|)$
        \ENDFOR
    \UNTIL{$\Delta < \theta$}
    \STATE \textbf{Output:} $V$ and greedy policy $\pi(s) = \argmax_{a \in \mc{A}(s)}\left(r(s,a) + \gamma\sum_{s' \in \mc{S}}p(s'|s,a)V(s')\right)$
    \end{algorithmic}
    \end{algorithm}

\section{Reinforcement Learning}

In section \ref{sec:dynamic_programming} a perfect model of the environment was used in order to compute the backup updates.
Model-Free Reinforcement Learning aims to solve the prediction and control problem only through experience borrowing ideas from Dynamic Programming.


\subsection{Temporal-Difference Learning}
\label{subsec:TD}
Temporal Difference Learning (TD) is a \textbf{bootstrapping} method: it learns directly from transitions of incomplete episodes using previous estimates.

\subsubsection{Prediction}
For the prediction problem TD updates its estimate of $v^\pi(s_t)$ only using the next step reward $r_{t+1}$ and its current estimate of $v^\pi(s_{t+1})$.

\begin{equation}
    \label{eq:TD(0)}
    V(S_t) \gets V(S_{t+1}) + \alpha [R_{t+1} + \gamma V(S_{t+1}) - V(S_t)]
\end{equation}

$R_{t+1} + \gamma V(S_{t+1})$ is called TD \textit{target} and the quantity $\delta_t \doteq R_{t+1} + \gamma V(S_{t+1}) - V(S_t)$ is called TD \textit{error}.



\subsubsection{Control}

For the Control problem TD algorithms can be divided into two categories:
\begin{enumerate}
    \item On-Policy Learning
    \item Off-Policy Learning
\end{enumerate}

In On Policy Learning the algorithmy directly improves the policy $\pi$ used to collect the samples from the environment while
Off Policy Learning allows to improve a policy $\pi'$ different from the one that interacts with the environment.
\\
One major problem in Model-Free reinforcement learning is the \textit{Exploration-Exploitation} tradeoff:
Should the agent try out new actions to discover their potential rewards or exploit its current knowledge in order to maximize the expected reward?
\\
\begin{definition}
    A soft policy $\pi$ is a policy with $\pi(a|s) > 0$, $\forall s \in \mc{S}$ and $\forall a \in \mc{A}(s)$.
\end{definition}

\begin{definition}
    An $\epsilon$-soft policy $\pi$ is a policy with $\pi(s,a) \ge \frac{\epsilon}{|\mc{A}(s)|}$, $\forall s \in \mc{S}$ and $\forall a \in \mc{A}(s)$ and some $\epsilon > 0$.
\end{definition}

\begin{definition}
    An $\epsilon$-greedy policy $\pi$ is a policy with:
    \begin{equation*}
        \pi(a|s) =
        \begin{cases}
            \frac{\epsilon}{|\mc{A}|} + 1 - \epsilon & \text{if } a = \argmax_{a' \in \mc{A}}Q(s, a') \\
            \frac{\epsilon}{|\mc{A}|}                & \text{otherwise}
        \end{cases}
    \end{equation*}
\end{definition}

The $\epsilon$-greedy policy is a specific example of $\epsilon$-soft policy that is the closest to a greedy policy.
It is the simplest idea for ensuring continual exploration while acting greedily.

\begin{theorem}
    For any $\epsilon$-greedy policy $\pi$, the $\epsilon$-greedy policy $\pi'$ with respect to $q^\pi$ is an improvement.
    \begin{align*}
        q^\pi(s,\pi'(s))    &= \sum_{a \in \mc{A}}\pi'(a|s)q^\pi(s,a) \\
                            &= \frac{\epsilon}{m}\sum_{a \in \mc{A}}q^\pi(s,a) + (1 - \epsilon)\max_{a \in \mc{A}}q^\pi(s,a) \\
                            &\ge \frac{\epsilon}{m}\sum_{a \in \mc{A}}q^\pi(s,a) + (1 - \epsilon)\sum_{a \in \mc{A}}\frac{\pi(a|s) - \frac{\epsilon}{m}}{1 - \epsilon}q^\pi(s,a) \\
                            &= \sum_{a \in \mc{A}}\pi(a|s)q^\pi(s,a) = v^\pi(s)
    \end{align*}
    with $m = |\mc{A}|$. \\
    Therefore from the policy improvement theorem \ref{th:policy_improvement}, $v^{\pi'}(s) \ge v^\pi(s)$.
\end{theorem}

\subsubsection{Sarsa: On-policy TD Control}
Applying the same TD method \eqref{eq:TD(0)} used to estimate $v^\pi$ to the action-value function $q^\pi$ leads to the following update rule:
\begin{equation*}
    Q(S_t, A_t) \gets Q(S_t, A_t) + \alpha [R_{t+1} + \gamma Q(S_{t+1}, A_{t+1}) - Q(S_t, A_t)]
\end{equation*}.

This rule uses every element of a quintuple of events, $(S_t, A_t, R_{t+1}, S_{t+1}, A_{t+1})$, that make up a transition from one state-action pair to the next.
This quintuple gives rise to the name \textit{Sarsa} for the algorithm.

It is straightforward to design an on-policy control algorithm based on the Sarsa prediction method. As in all on-policy methods, we continually estimate $q^\pi$
for the behavior policy $\pi$, and at the same time change $\pi$ towards greediness with respect to $q^\pi$.

\begin{algorithm}[H]
    \caption{Sarsa}
    \label{alg:sarsa}
    \begin{algorithmic}[1]
        \STATE Initialize $Q(s,a)$ arbitrarily
        \LOOP
            \STATE Initialize $S$
            \STATE Choose $A$ from $S$ using policy derived from $Q$ (e.g. $\epsilon$-greedy)
            \REPEAT
                \STATE Take action $A$, observe $R$,$S'$
                \STATE Choose $A'$ from $S'$ using policy derived from $Q$ (e.g. $\epsilon$-greedy)
                \STATE $Q(S,A) \gets Q(S,A) + \alpha[R + \gamma Q(S',A') - Q(S, A)]$
                \STATE $S \gets S'$; $A \gets A'$
            \UNTIL{$S$ is terminal}
        \ENDLOOP
    \end{algorithmic}
\end{algorithm}

\subsubsection{Q-Learning: Off-policy TD Control}
One of the early breakthroughs in reinforcement learning was the development of an off-policy TD control algorithm known as \textit{Q-Learning}
\cite{Watkins:1989}.

In this algorithm the learned action-value function $Q$ directly approximates $q^*$, the optimal action-value function, independent of the 
policy being followed:
\begin{equation*}
    Q(S_t,A_t) \gets Q(S_t,A_t) + \alpha [R_{t+1} + \gamma \max_{a \in \mc{A}}Q(S_{t+1}, a) - Q(S_t, A_t)]
\end{equation*}

In this case, the learned action-value function $Q$ directly approximates  $q^*$, the optimal value action-value function thanks to the $\max$ operation.

\begin{algorithm}[H]
    \caption{Q-Learning}
    \label{alg:q-learning}
    \begin{algorithmic}[1]
        \STATE Initialize $Q(s,a)$ arbitrarily
        \LOOP
            \STATE Initialize $S$
            \STATE Choose $A$ from $S$ using policy derived from $Q$ (e.g. $\epsilon$-greedy)
            \REPEAT
                \STATE Take action $A$, observe $R$,$S'$
                \STATE $Q(S,A) \gets Q(S,A) + \alpha[R + \gamma \max_{a \in \mc{A}}Q(S',a) - Q(S, A)]$
                \STATE $S \gets S'$
            \UNTIL{$S$ is terminal}
        \ENDLOOP
    \end{algorithmic}
\end{algorithm}

\subsubsection{Convergence of Sarsa and Q-Learning}
\begin{definition}
    A learning policy is called Greedy in the limit of infinite exploration (GLIE) if it satisfies the following two conditions:
    \begin{enumerate}
        \item All state-action pairs are explored infinitely many times
        \begin{equation*}
            \lim_{k \to \infty}N_k(s,a) = \infty
        \end{equation*}
        \item The policy converges on a greedy policy
        \begin{equation*}
            \lim_{k \to \infty} \pi_k(a|s) = \mathbbm{1}(a = \argmax_{a' \in \mc{A}}Q_k(s,a'))
        \end{equation*}
    \end{enumerate}
\end{definition}

\begin{theorem}
    SARSA converges to the optimal action-value function $Q(s,a) \to Q^*(s,a)$ under the following conditions:
    \begin{enumerate}
        \item GLIE sequence of policies $\pi_t(s,a)$
        \item Robbins-Monro sequence of step-sizes $\alpha_t$
            \begin{align*}
                &\sum_{t=1}^\infty \alpha_t = \infty
                &\sum_{t=1}^\infty \alpha_t^2 < \infty
            \end{align*}
    \end{enumerate}
\end{theorem}

This means that in order for SARSA to converge to $q^*$ we need to decrease $\epsilon$ over time (e.g. $\epsilon = \frac{1}{t}$).
If $\epsilon$ is kept fixed SARSA will converge to the optimal policy among the $\epsilon$-soft policies.
Being Q-Learning an off-policy algorithm that is already improving Q towards $q^*$ there is no need for an $\epsilon$ decay as long as all state-action
pairs are visited infinitely and the learning rate' schedule satisfies the Robbins-Monro conditions.


\subsection{Reinforcement Learning Taxonomy}
% Tutte le distinzioni dei vari algoritmi
This subsection includes a brief overview of the whole Reinforcement Learning taxonomy. \ref{subsec:TD} presented two well known TD tabular algorithms
but often real world scenarios have to be modeled with complex state and action spaces with more dimensions resulting in intractable problems.
%Moreover these spaces could be continuous and tabular methods cannot be applied anymore.

\begin{itemize}
\item \textbf{Function Approximation} \\
When the state space is too big for tabular reinforcement learning approximate solution methods are used instead.
Typically the value function (or action-value function) is parameterized with a weight vector $\mathbf{w} \in \mathbb{R}^d$ and $d \ll |\mc{S}|$:
\begin{equation*}
    \hat v(s, \mathbf{w}) \approx v_\pi(s)
\end{equation*}

For example $\hat v$ might be a linear function in features of the state, with $\mathbf{w}$ the vector of feature weights.
More generally, $\hat v$ might be the function computed by a multi-layer artificial neural network or even the function computed
by a decision tree.
\\

In addition, function approximation allow continous state spaces to be managed without the need for discretization techniques.

\item \textbf{Value-Based Algorithms} \\
All the algorithms that seek to learn the action values and then select actions based on their estimated values are called Value-Based Algorithms.
The temporal difference algorithms presented in \ref{subsec:TD} are all value-based algorithms.
\item \textbf{Policy-Based Algorithms} \\
Differently from Value-Based Algorithms Policy-Based algorithms seek to learn a \textit{parameterized policy} that can select actions without consulting
a value function:
\begin{equation*}
    \pi(a|s,\theta) = \mathbb{P}(A_t = a | S_t = s, \theta_t = \theta)
\end{equation*}

One important consequence of the use of parameterized policies is that they offer practical ways of dealing with large action spaces,
even continuous spaces with an infinite number of actions. Instead of computing the learned probability for each of the many actions,
we instead learn statistics of the probability distribution. For example, the action set might be the real numbers, with actions chosen
from a normal distribution:

\begin{definition}
    A gaussian policy suitable for continuous action space is given by
    \begin{equation}
        \label{eq:gaussian_policy}
        \pi(a|s,\theta) = \frac{1}{\sigma(s, \theta)\sqrt{2 \pi}}\exp\left(-\frac{(a - \mu(s, \theta))^2}{2\sigma(s,\theta)^2}\right)
    \end{equation}
\end{definition}

where $\mu: \mc{S} \times \Theta \to \mathbb{R}$ is the mean function and $\sigma: \mc{S} \times \Theta \to (0, \infty)$ is the standard deviation function.
$\Theta$ is the parameter space of the policy.

A sample from a gaussian policy can be computed as follows:
\begin{equation*}
    a = \mu(s, \theta) + \sigma(s, \theta) \eta
\end{equation*}

drawing $\eta$ from a standard distribution $\mc{N}(0,1)$.



\item \textbf{Actor-Critic Algorithms} \\
A combination of value-based methods and policy-based methods gives rise to Actor-Critic algorithms. Actor critic algorithms
learn a parameterized policy (Actor) and a value function (Critic).
The actor is updated in the direction that maximizes the expected reward using gradients derived from the critic's feedback.
The critic estimates the actions taken by the actor.

\end{itemize}

\subsection{Deep Reinforcement Learning}
In recent years Deep Learning techniques applied to Reinforcement Learning techniques have shown to yield very good results.
The first notable example came from Mnih et al. \cite{DQN} modifying Q-Learning to work with Deep Neural Networks.
At that time, their approach achieved state of the art on multiple Atari Games with no prior knowledge and same hyperparameters for each game. This algorithm, called
\textit{Deep Q-Networks} (DQN) uses an \textit{experience replay buffer} to store played transitions for sampling minibatches.
These minibatches are used to train the Q-Network using Stochastic Gradient Descent with backpropagation on the neural network.

Unfortunately nonlinear function approximators like neural networks come with a theoretical cost:
convergence to an optimum is only guaranteed locally. Despite this limitation, deep reinforcement learning approaches remain the most widely used methods in the field 
due to their empirical success and ability to handle complex, high-dimensional state spaces.


Below, a state of the art algorithm called \textit{Soft Actor-Critic} (SAC) is presented.

\subsubsection{Soft Actor-Critic}
SAC \cite{SAC, SAC_modified} is an off policy actor-critic deep RL algorithm based on the maximum entropy framework in which the standard reinforcement learning
objective is augmented with an entropy term:
\begin{equation*}
    \pi^* = \argmax_\pi \sum_t \mathbb{E}_{(s_t, a_t) \sim \rho_\pi}[r(s_t, a_t) + \alpha \mc{H}(\pi(.|s_t))]
\end{equation*}

where $\alpha$ is the temperature parameter that determines the relative importance of the entropy term versus the reward, and thus controls the stochasticity
of the optimal policy.

This objective has shown to improve exploration and improve learning speed over state-of-art methods that optimize the conventional RL objective function.

\begin{definition}
    The soft state-value and action value functions  $v_{\text{soft}}: \mc{S} \to \mathbb{R}$ and $q_{\text{soft}}: \mc{S} \times \mc{A} \to \mathbb{R}$ are defined by:
    \begin{align*} 
        &v_\text{soft}^\pi(s_t) = \mathbb{E}_{a_t \sim \pi }[q^\pi(s_t, a_t) - \alpha \log \pi(a_t | s_t)] \\
        &q_\text{soft}^\pi(s_t, a_t) = r(s_t, a_t) + \gamma \mathbb{E}_{s_{t+1} \sim \mathit{p}}[v_\text{soft}^\pi(s_{t+1})]
    \end{align*}
\end{definition}

\textit{Soft Actor-Critic} is based on the soft policy iteration algorithm, an adaptation of policy iteration algorithm that includes the entropy term in the value function
and guarantees convergence to an optimal policy.

In the soft policy iteration, the policy is updated towards the exponential of the soft Q-function. This particular choice of update can be guaranteed
to result in an improved policy in terms of its soft values:

\begin{equation}
    \label{eq:soft_policy_improvement}
    \pi_\text{new} = \argmin_{\pi' \in \Pi} D_{KL} \left(\pi'(.|s_t) \Big\Vert \frac{\exp(\frac{1}{\alpha}Q^\text{old}(s_t, .))}{Z^{\pi_{\text{old}}(s_t)}}\right),
\end{equation}

where the partition function $Z^\pi_\text{old}(s_t)$ normalizes the distribution, and while it is intractable in general, it does not contribute to the gradient with respect to the new policy and thus can be ignored.

\textit{Soft Actor-Critic} uses function approximators for both the soft Q-function and the policy, and instead of running evaluation and improvement to convergence
like soft policy iteration, alternates between optimizing both networks with stochastic gradient descent. The soft Q-function can be modeled
as expressive neural networks, and the policy as a Gaussian distribution \eqref{eq:gaussian_policy} with mean and covariance given by neural networks 

The soft Q-function parameters can be trained to minimize the soft Bellman residual through stochastic gradients:

\begin{equation*}
    J_Q(\theta) = \mathbb{E}_{(s_t,a_t) \sim \mathcal{D}} \left[ \frac{1}{2} \left( Q_\theta(s_t, a_t) - \left( r(s_t, a_t) + \gamma \mathbb{E}_{s_{t+1} \sim p} [V_{\bar{\theta}}(s_{t+1})] \right) \right)^2 \right]
\end{equation*}

where $\mc{D}$ is the replay buffer storing past experience. In whis way the reinforcement learning objective is approximated
as a supervised learning objective over a continually changing buffer of past experience.

\begin{equation*}
    \hat{\nabla}_\theta J_Q(\theta) = \nabla_\theta Q_\theta(a_t, s_t) \left( Q_\theta(s_t, a_t) - \left( r(s_t, a_t) + \gamma \left( Q_{\bar{\theta}}(s_{t+1}, a_{t+1}) - \alpha \log (\pi(a_{t+1} | s_{t+1})) \right) \right) \right)
\end{equation*}

The update makes use of a target Q-Function with parameters $\bar{\theta}$ that are obtained as an exponential moving average of the soft Q-function weights, which has shown to stabilize training \cite{DQN}.
Finally, the policy parameters can be learned by directly minimizing the expected KL-Divergence in \eqref{eq:soft_policy_improvement}:

\begin{equation}
    \label{eq:sac_actor_loss}
    J_\pi(\phi) = \mathbb{E}_{s_t \sim \mc{D}} \left[\mathbb{E}_{a_t \sim \pi_\phi} \left[\alpha \log{\pi_\phi(a_t|s_t)} - Q_\theta(s_t, a_t)\right]\right]
\end{equation}

As the Q-function is differentiable it is convenient to apply the reparameterization trick to the policy in order to minimize $J_\pi$:

\begin{equation*}
    a_t = f_\phi(\epsilon_t; s_t)
\end{equation*}

where $\epsilon_t$ is an input noise, sampled from some fixed distribution, such as a spherical Gaussian.

The objective function in \eqref{eq:sac_actor_loss} can be rewritten as
\begin{equation}
    \label{eq:sac_actor_loss2}
    J_\pi(\phi) = \mathbb{E}_{s_t \sim \mc{D}, \epsilon_t \sim \mc{N}}\left[\alpha \log{\pi_\phi(f_\phi(\epsilon_t;s_t)|s_t)} - Q_\theta(s_t, f_\phi(\epsilon_t;s_t))\right]
\end{equation}

where $\pi_\phi$ is implicitly defined in terms of $f_\phi$. The gradient of equation \eqref{eq:sac_actor_loss2} is then approximated with

\begin{equation*}
    \hat{\nabla}_\phi J_\pi(\phi) = \nabla_\phi \alpha \log{\pi_\phi(a_t|s_t)} + \left(\nabla_{a_t} \alpha \log{\pi_\phi(a_t|s_t)} - \nabla_{a_t}Q_\theta(s_t, a_t)\right) \nabla_\phi f_\phi(\epsilon_t;s_t)
\end{equation*}
where $a_t$ is evaluated at $f_\phi(\epsilon_t;s_t)$.

SAC also employs an automatic entropy adjustment algorithm that finds the optimal entropy coefficient during training.
Also, it makes use of two soft Q functions to mitigate positive bias in the policy improvement step that is known to degrade performance of value based methods \cite{Hasselt2010, Fujimoto2018}.
In particular, we parameterize two soft Q-functions, with parameters $\theta_i$ , and train them independently to optimize $J_Q(\theta_i)$.
For further details refer to \cite{SAC_modified}.



\begin{algorithm}[H]
    \caption{Soft Actor-Critic (SAC) Algorithm}
    \begin{algorithmic}[1]
    \STATE \textbf{Initialize} policy parameters $\phi$, Q-function parameters $\theta_1$, $\theta_2$
    \STATE \textbf{Initialize} replay buffer $\mathcal{D}$
    \STATE \textbf{Set} target Q-function weights: $\bar{\theta}_i \leftarrow \theta_i$ for $i \in \{1, 2\}$
    \FOR{each iteration}
        \FOR{each environment step}
            \STATE Sample action $a_t \sim \pi_\theta(a_t|s_t)$
            \STATE Execute action $a_t$ in the environment
            \STATE Observe next state $s_{t+1}$, reward $r_t$
            \STATE Store $(s_t, a_t, r_t, s_{t+1})$ in replay buffer $\mathcal{D}$
        \ENDFOR
        \FOR{each gradient step}
            \STATE Sample a batch of transitions $(s, a, r, s')$ from $\mathcal{D}$
            \STATE $\theta_i \leftarrow \theta_i - \lambda_Q \hat{\nabla}_{\theta_i}J_Q(\theta_i)$ for $i \in \{1, 2\}$
            \STATE $\phi \leftarrow \phi - \lambda_\pi \hat{\nabla}_\phi J_\pi(\phi)$
            \STATE $\bar{\theta}_i \leftarrow \tau \theta_i + (1 - \tau) \bar{\theta}_i$ for $i \in \{1, 2\}$ 
            
        \ENDFOR
    \ENDFOR
    \end{algorithmic}
    \end{algorithm}
















% \chapter{Chapter one}
\label{ch:chapter_one}%
% The \label{...}% enables to remove the small indentation that is generated, always leave the % symbol.

In this chapter additional useful information are reported.

\section{Sections and subsections}
\label{sec:section_name}
Chapters are typically subdivided into sections and subsections, and, optionally,
subsubsections, paragraphs and subparagraphs.
All can have a title, but only sections and subsections are numbered.
A new section is created by the command
\begin{verbatim}
\section{Title of the section}
\end{verbatim}
The numbering can be turned off by using \verb|\section*{}|.
\\
A new subsection is created by the command
\begin{verbatim}
\subsection{Title of the subsection}
\end{verbatim}
and, similarly, the numbering can be turned off by adding an asterisk as follows 
\begin{verbatim}
\subsection*{}
\end{verbatim}

\section{Equations}
\label{sec:eqs}
This section gives some examples of writing mathematical equations in your thesis.

Maxwell's equations read:
\begin{subequations}
    \label{eq:maxwell}
    \begin{align}[left=\empheqlbrace]
    \nabla\cdot \bm{D} & = \rho, \label{eq:maxwell1} \\
    \nabla \times \bm{E} +  \frac{\partial \bm{B}}{\partial t} & = \bm{0}, \label{eq:maxwell2} \\
    \nabla\cdot \bm{B} & = 0, \label{eq:maxwell3} \\
    \nabla \times \bm{H} - \frac{\partial \bm{D}}{\partial t} &= \bm{J}. \label{eq:maxwell4}
    \end{align}
\end{subequations}

Equation~\eqref{eq:maxwell} is automatically labeled by \texttt{cleveref},
as well as Equation~\eqref{eq:maxwell1} and Equation~\eqref{eq:maxwell3}.
Thanks to the \verb|cleveref| package, there is no need to use \verb|\eqref|.
Remember that Equations have to be numbered only if they are referenced in the text.

Equations~\eqref{eq:maxwell_multilabels1}, \eqref{eq:maxwell_multilabels2}, \eqref{eq:maxwell_multilabels3}, and \eqref{eq:maxwell_multilabels4} show again Maxwell's equations without brace:
\begin{align}
    \nabla\cdot \bm{D} & = \rho, \label{eq:maxwell_multilabels1} \\
    \nabla \times \bm{E} +  \frac{\partial \bm{B}}{\partial t} &= \bm{0}, \label{eq:maxwell_multilabels2} \\
    \nabla\cdot \bm{B} & = 0, \label{eq:maxwell_multilabels3} \\
    \nabla \times \bm{H} - \frac{\partial \bm{D}}{\partial t} &= \bm{J} \label{eq:maxwell_multilabels4}.
\end{align}

Equation~\eqref{eq:maxwell_singlelabel} is the same as before,
but with just one label:
\begin{equation}
    \label{eq:maxwell_singlelabel}
    \left\{
    \begin{aligned}
    \nabla\cdot \bm{D} & = \rho, \\
    \nabla \times \bm{E} +  \frac{\partial \bm{B}}{\partial t} &= \bm{0},\\
    \nabla\cdot \bm{B} & = 0, \\
    \nabla \times \bm{H} - \frac{\partial \bm{D}}{\partial t} &= \bm{J}.
    \end{aligned}
    \right.
\end{equation}

\section{Figures, Tables and Algorithms}
Figures, Tables and Algorithms have to contain a Caption that describe their content, and have to be properly reffered in the text.

\subsection{Figures}
\label{subsec:figures}

For including pictures in your text you can use \texttt{TikZ} for high-quality hand-made figures,
or just include them as usual with the command
\begin{verbatim}
\includegraphics[options]{filename.xxx}
\end{verbatim}
Here xxx is the correct format, e.g. \verb|.png|, \verb|.jpg|, \verb|.eps|, \dots.

\begin{figure}[H]
    \centering
    \includegraphics[width=0.3\textwidth]{logo_polimi_scritta.eps}
    \caption{Caption of the Figure to appear in the List of Figures.}
    \label{fig:quadtree}
\end{figure}

Thanks to the \texttt{\textbackslash subfloat} command, a single figure, such as Figure~\ref{fig:quadtree},
can contain multiple sub-figures with their own caption and label, e.g. \color{black} Figure~\ref{fig:polimi_logo1} and Figure~\ref{fig:polimi_logo2}. 

\begin{figure}[H]
    \centering
    \subfloat[One PoliMi logo.\label{fig:polimi_logo1}]{
        \includegraphics[scale=0.5]{Images/logo_polimi_scritta.eps}
    }
    \quad
    \subfloat[Another one PoliMi logo.\label{fig:polimi_logo2}]{
        \includegraphics[scale=0.5]{Images/logo_polimi_scritta2.eps}
    }
    \caption[Shorter caption]{This is a very long caption you don't want to appear in the List of Figures.}
    \label{fig:quadtree2}
\end{figure}


\subsection{Tables}
\label{subsec:tables}

Within the environments \texttt{table} and  \texttt{tabular} you can create very fancy tables as the one shown in Table~\ref{table:example}.
\begin{table}[H]
    \caption*{\textbf{Title of Table (optional)}}
    \centering 
    \begin{tabular}{|p{3em} c c c |}
    \hline
    \rowcolor{bluepoli!40} % comment this line to remove the color
     & \textbf{column 1} & \textbf{column 2} & \textbf{column 3} \T\B \\
    \hline \hline
    \textbf{row 1} & 1 & 2 & 3 \T\B \\
    \textbf{row 2} & $\alpha$ & $\beta$ & $\gamma$ \T\B\\
    \textbf{row 3} & alpha & beta & gamma \B\\
    \hline
    \end{tabular}
    \\[10pt]
    \caption{Caption of the Table to appear in the List of Tables.}
    \label{table:example}
\end{table}

You can also consider to highlight selected columns or rows in order to make tables more readable.
Moreover, with the use of \texttt{table*} and the option \texttt{bp} it is possible to align them at the bottom of the page. One example is presented in Table~\ref{table:exampleC}. 

\begin{table}[H]
\centering 
    \begin{tabular}{|p{3em} | c | c | c | c | c | c|}
    \hline
%    \rowcolor{bluepoli!40}
     & \textbf{column1} & \textbf{column2} & \textbf{column3} & \textbf{column4} & \textbf{column5} & \textbf{column6} \T\B \\
    \hline \hline
    \textbf{row1} & 1 & 2 & 3 & 4 & 5 & 6 \T\B\\
    \textbf{row2} & a & b & c & d & e & f \T\B\\
    \textbf{row3} & $\alpha$ & $\beta$ & $\gamma$ & $\delta$ & $\phi$ & $\omega$ \T\B\\
    \textbf{row4} & alpha & beta & gamma & delta & phi & omega \B\\
    \hline
    \end{tabular}
    \\[10pt]
    \caption{Highlighting the columns}
    \label{table:exampleC}
\end{table}

\begin{table}[H]
\centering 
    \begin{tabular}{|p{3em} c c c c c c|}
    \hline
%    \rowcolor{bluepoli!40}
     & \textbf{column1} & \textbf{column2} & \textbf{column3} & \textbf{column4} & \textbf{column5} & \textbf{column6} \T\B \\
    \hline \hline
    \textbf{row1} & 1 & 2 & 3 & 4 & 5 & 6 \T\B\\
    \hline
    \textbf{row2} & a & b & c & d & e & f \T\B\\
    \hline
    \textbf{row3} & $\alpha$ & $\beta$ & $\gamma$ & $\delta$ & $\phi$ & $\omega$ \T\B\\
    \hline
    \textbf{row4} & alpha & beta & gamma & delta & phi & omega \B\\
    \hline
    \end{tabular}
    \\[10pt]
    \caption{Highlighting the rows}
    \label{table:exampleR}
\end{table}

\subsection{Algorithms}
\label{subsec:algorithms}

Pseudo-algorithms can be written in \LaTeX{} with the \texttt{algorithm} and \texttt{algorithmic} packages.
An example is shown in Algorithm~\ref{alg:var}.
\begin{algorithm}[H]
    \label{alg:example}
    \caption{Name of the Algorithm}
    \label{alg:var}
    \label{protocol1}
    \begin{algorithmic}[1]
    \STATE Initial instructions
    \FOR{$for-condition$}
    \STATE{Some instructions}
    \IF{$if-condition$}
    \STATE{Some other instructions}
    \ENDIF
    \ENDFOR
    \WHILE{$while-condition$}
    \STATE{Some further instructions}
    \ENDWHILE
    \STATE Final instructions
    \end{algorithmic}
\end{algorithm} 

\vspace{5mm}

\section{Theorems, propositions and lists}

\subsection{Theorems}
Theorems have to be formatted as:
\begin{theorem}
\label{a_theorem}
Write here your theorem. 
\end{theorem}
\textit{Proof.} If useful you can report here the proof.

\subsection{Propositions}
Propositions have to be formatted as:
\begin{proposition}
Write here your proposition.
\end{proposition}

\subsection{Lists}
How to  insert itemized lists:
\begin{itemize}
    \item first item;
    \item second item.
\end{itemize}
How to insert numbered lists:
\begin{enumerate}
    \item first item;
    \item second item.
\end{enumerate}

\section{Use of copyrighted material}

Each student is responsible for obtaining copyright permissions, if necessary, to include published material in the thesis.
This applies typically to third-party material published by someone else.

\section{Plagiarism}

You have to be sure to respect the rules on Copyright and avoid an involuntary plagiarism.
It is allowed to take other persons' ideas only if the author and his original work are clearly mentioned.
As stated in the Code of Ethics and Conduct, Politecnico di Milano \textit{promotes the integrity of research,
condemns manipulation and the infringement of intellectual property}, and gives opportunity to all those
who carry out research activities to have an adequate training on ethical conduct and integrity while doing research.
To be sure to respect the copyright rules, read the guides on Copyright legislation and citation styles available
at:
\begin{verbatim}
https://www.biblio.polimi.it/en/tools/courses-and-tutorials
\end{verbatim}
You can also attend the courses which are periodically organized on "Bibliographic citations and bibliography management".

\section{Bibliography and citations}
Your thesis must contain a suitable Bibliography which lists all the sources consulted on developing the work.
The list of references is placed at the end of the manuscript after the chapter containing the conclusions.
We suggest to use the BibTeX package and save the bibliographic references  in the file \verb|Thesis_bibliography.bib|.
This is indeed a database containing all the information about the references. To cite in your manuscript, use the \verb|\cite{}| command as follows:
\\
\textit{Here is how you cite bibliography entries: \cite{knuth74}, or multiple ones at once: \cite{knuth92,lamport94}}.
\\
The bibliography and list of references are generated automatically by running BibTeX \cite{bibtex}.

\chapter{Conclusions and future developments}
\label{ch:conclusions}%
In this work, reinforcement learning was successfully applied to the air hockey task using a general-purpose manipulator. 
The agent was able to play entire matches without violating too many constraints, thanks to the ATACOM algorithm. 
The SAC method outperformed the rule-based controller for the hitting task, achieving a higher success rate. 
The rule-based switcher effectively selected the appropriate policy at the right time, with its parameters slightly optimized PGPE.

\section{Future works}
The developed rule-based switcher currently uses a few simple parametric rules. 
Future work could involve implementing a more complex switcher with additional rules that also take the opponent's actions into account.

An other strategy could be to apply an end-to-end hierarchical reinforcement learning approach, such as HIRO \cite{HIRO}, where a single low-level policy receives from the high-level
policy a vector to determine the desired behavior.

Since air hockey is a multi-agent setting, investigating a self-play learning strategy could also be beneficial. 
This approach could help the agent improve by continuously adapting to and learning from interactions with itself or other agents.



% %-------------------------------------------------------------------------
%	BIBLIOGRAPHY
%-------------------------------------------------------------------------

\addtocontents{toc}{\vspace{2em}} % Add a gap in the Contents, for aesthetics
\bibliography{Thesis_bibliography} % The references information are stored in the file named "Thesis_bibliography.bib"

%-------------------------------------------------------------------------
%	BIBLIOGRAPHY
%-------------------------------------------------------------------------

\addtocontents{toc}{\vspace{2em}} % Add a gap in the Contents, for aesthetics
\bibliography{Thesis_bibliography} % The references information are stored in the file named "Thesis_bibliography.bib"


%-------------------------------------------------------------------------
%	APPENDICES
%-------------------------------------------------------------------------

\cleardoublepage
\addtocontents{toc}{\vspace{2em}} % Add a gap in the Contents, for aesthetics
\appendix

\chapter{Appendix A}
If you need to include an appendix to support the research in your thesis, you can place it at the end of the manuscript.
An appendix contains supplementary material (figures, tables, data, codes, mathematical proofs, surveys, \dots)
which supplement the main results contained in the previous chapters.

\chapter{Appendix B}
It may be necessary to include another appendix to better organize the presentation of supplementary material.



% LIST OF FIGURES
\listoffigures

% LIST OF TABLES
\listoftables

% LIST OF SYMBOLS
% Write out the List of Symbols in this page
% \chapter*{List of Symbols} % You have to include a chapter for your list of symbols (
% \begin{table}[H]
%     \centering
%     \begin{tabular}{lll}
%         \textbf{Variable} & \textbf{Description} & \textbf{SI unit} \\\hline\\[-9px]
%         $\bm{u}$ & solid displacement & m \\[2px]
%         $\bm{u}_f$ & fluid displacement & m \\[2px]
%     \end{tabular}
% \end{table}

% ACKNOWLEDGEMENTS
% \chapter*{Acknowledgements}
% Here you might want to acknowledge someone.


\cleardoublepage


\end{document}
